\documentclass[10pt,a4paper]{amsart}
\usepackage{amsmath,amssymb,amsthm,graphicx}
\usepackage{verbatim}
\newenvironment{reference}{\comment}{\endcomment}
%\newenvironment{reference}{}{}
\newenvironment{slogan}{\comment}{\endcomment}
\newenvironment{history}{\comment}{\endcomment}

% For commutative diagrams we use Xy-pic
\usepackage[all]{xy}

% We use 2cell for 2-commutative diagrams.
\xyoption{2cell}
% Theorem environments.
%
\theoremstyle{plain}
\newtheorem{theorem}[subsection]{Theorem}
\newtheorem{proposition}[subsection]{Proposition}
\newtheorem{lemma}[subsection]{Lemma}

\theoremstyle{definition}
\newtheorem{definition}[subsection]{Definition}
\newtheorem{example}[subsection]{Example}
\newtheorem{exercise}[subsection]{Exercise}
\newtheorem{situation}[subsection]{Situation}

\theoremstyle{remark}
\newtheorem{remark}[subsection]{Remark}
\newtheorem{remarks}[subsection]{Remarks}

\numberwithin{equation}{subsection}

% Macros
%
\def\lim{\mathop{\mathrm{lim}}\nolimits}
\def\colim{\mathop{\mathrm{colim}}\nolimits}
\def\Spec{\mathop{\mathrm{Spec}}}
\def\Hom{\mathop{\mathrm{Hom}}\nolimits}
\def\Ext{\mathop{\mathrm{Ext}}\nolimits}
\def\SheafHom{\mathop{\mathcal{H}\!\mathit{om}}\nolimits}
\def\SheafExt{\mathop{\mathcal{E}\!\mathit{xt}}\nolimits}
\def\Sch{\mathit{Sch}}
\def\Mor{\mathop{\mathrm{Mor}}\nolimits}
\def\Ob{\mathop{\mathrm{Ob}}\nolimits}
\def\Sh{\mathop{\mathit{Sh}}\nolimits}
\def\NL{\mathop{N\!L}\nolimits}
\def\CH{\mathop{\mathrm{CH}}\nolimits}
\def\proetale{{pro\text{-}\acute{e}tale}}
\def\etale{{\acute{e}tale}}
\def\QCoh{\mathit{QCoh}}
\def\Ker{\mathop{\mathrm{Ker}}}
\def\Im{\mathop{\mathrm{Im}}}
\def\Coker{\mathop{\mathrm{Coker}}}
\def\Coim{\mathop{\mathrm{Coim}}}

% Boxtimes
%
\DeclareMathSymbol{\boxtimes}{\mathbin}{AMSa}{"02}

%
% Macros for moduli stacks/spaces
%
\def\QCohstack{\mathcal{QC}\!\mathit{oh}}
\def\Cohstack{\mathcal{C}\!\mathit{oh}}
\def\Spacesstack{\mathcal{S}\!\mathit{paces}}
\def\Quotfunctor{\mathrm{Quot}}
\def\Hilbfunctor{\mathrm{Hilb}}
\def\Curvesstack{\mathcal{C}\!\mathit{urves}}
\def\Polarizedstack{\mathcal{P}\!\mathit{olarized}}
\def\Complexesstack{\mathcal{C}\!\mathit{omplexes}}
% \Pic is the operator that assigns to X its picard group, usage \Pic(X)
% \Picardstack_{X/B} denotes the Picard stack of X over B
% \Picardfunctor_{X/B} denotes the Picard functor of X over B
\def\Pic{\mathop{\mathrm{Pic}}\nolimits}
\def\Picardstack{\mathcal{P}\!\mathit{ic}}
\def\Picardfunctor{\mathrm{Pic}}
\def\Deformationcategory{\mathcal{D}\!\mathit{ef}}
\usepackage{cleveref}
\crefname{originalreference}{original reference}{original references}
\Crefname{originalreference}{Original reference}{Original references}
\makeatletter
\newlabel{lemma-contract-ample}{{0C2M}{source}{}{sourcepointer-1}{}}
\newlabel{lemma-contract-ample@cref}{{[originalreference][0][]0C2M}{[1][0][]source}{}{}{}}
\makeatother
\crefname{originaldefinition}{original definition}{original definitions}
\Crefname{originaldefinition}{Original definition}{Original definitions}
\makeatletter
\newlabel{morphisms-definition-quasi-projective}{{01VW}{source}{}{sourcepointer-2}{}}
\newlabel{morphisms-definition-quasi-projective@cref}{{[originaldefinition][0][]01VW}{[1][0][]source}{}{}{}}
\makeatother
\crefname{originallemma}{original lemma}{original lemmas}
\Crefname{originallemma}{Original lemma}{Original lemmas}
\makeatletter
\newlabel{morphisms-lemma-characterize-relatively-ample}{{01VJ}{source}{}{sourcepointer-3}{}}
\newlabel{morphisms-lemma-characterize-relatively-ample@cref}{{[originallemma][0][]01VJ}{[1][0][]source}{}{}{}}
\makeatother
\crefname{originallemma}{original lemma}{original lemmas}
\Crefname{originallemma}{Original lemma}{Original lemmas}
\makeatletter
\newlabel{morphisms-lemma-projective-is-quasi-projective-proper}{{0BCL}{source}{}{sourcepointer-4}{}}
\newlabel{morphisms-lemma-projective-is-quasi-projective-proper@cref}{{[originallemma][0][]0BCL}{[1][0][]source}{}{}{}}
\makeatother
\crefname{originallemma}{original lemma}{original lemmas}
\Crefname{originallemma}{Original lemma}{Original lemmas}
\makeatletter
\newlabel{morphisms-lemma-proper-local-on-the-base}{{01W2}{source}{}{sourcepointer-5}{}}
\newlabel{morphisms-lemma-proper-local-on-the-base@cref}{{[originallemma][0][]01W2}{[1][0][]source}{}{}{}}
\makeatother
\crefname{originallemma}{original lemma}{original lemmas}
\Crefname{originallemma}{Original lemma}{Original lemmas}
\makeatletter
\newlabel{morphisms-lemma-quasi-projective-open-projective}{{07RM}{source}{}{sourcepointer-6}{}}
\newlabel{morphisms-lemma-quasi-projective-open-projective@cref}{{[originallemma][0][]07RM}{[1][0][]source}{}{}{}}
\makeatother
\title{Contraction of an exceptional curve of the first kind preserving projectivity and quasi-projectivity}
\author{The Stacks Project Authors}
\date{stacks-project:resolve.tex:local-snapshot-20261012; GNU FDL 1.2 or later}
\begin{document}
\maketitle
\noindent\textit{Editorial note.} All delivered mathematical text is the author's own, copied from the pinned Stacks Project chapter file resolve.tex; no statement or proof step was rewritten, reordered or removed.\par\medskip
\noindent\textit{Editorial note.} The counted proof is the author's complete argument: the contraction over an affine open of the base together with the ample module produced there, the glueing of the contracted piece with the complement of the curve and the verification of properness by localness on the base, and the glueing of the invertible modules yielding relative ampleness followed by the reduction of the quasi-projective case to a projective model.\par\medskip
\noindent\textit{Editorial note.} The label lemma-contract-when-quasi-projective names the counted statement, and the proof environment immediately following it in the source is delivered verbatim as its proof.\par\medskip
\noindent\textit{Editorial note.} The counted statement is carried by the author's own lemma environment in the chapter source.\par\medskip
\noindent\textit{Editorial note.} Pre-open gate: python operations/agent-authoring/tools/precard\_gate.py --source resolve.tex --label lemma-contract-when-quasi-projective --start 4985 --end 5047 --local-id 28082 returned RESULT=OK (L28082 free, occupancy FREE, public interval clear). Pre-assembly tag self-check: canonical tag resolved and not present in reference/public-items.jsonl.\par\medskip
\noindent\textit{Editorial note.} The statement and proof are copied from the pinned chapter file; the cited inputs are kept as external pointers with their canonical Stacks tags.\par\medskip
\noindent\textit{Allowed standard prerequisites (editorial).}\begin{itemize}
\item Contraction over an affine base with an ample module on the target. Same chapter, the lemma producing a contraction of an exceptional curve of the first kind over an affine base together with an ample invertible module on the target whose pullback is a twist of the given module. Supplies the contraction over an affine open containing the image of the curve and the ample invertible module on the target.
\item Characterization of projectivity and of quasi-projectivity through relative ampleness. Morphisms, the lemmas relating a projective morphism to properness plus the existence of a relatively ample invertible module, and the lemma writing a quasi-projective scheme as an open subscheme of a projective one. Supply the description of the two hypotheses and the embedding used in the second part.
\item Localness of properness on the base. Morphisms, the lemma that properness can be checked after restricting to the members of an open covering of the base. Supplies the conclusion that the glued morphism is proper over the base.
\end{itemize}\medskip
\ifdefined\hypertarget\hypertarget{sourcepointer-1}{}\fi
\noindent\textit{Source reference.} Original reference 0C2M: sources/stacks-resolve/pinned-original-source.tex, source lines 4883--4883. Prerequisite pointer invoked at proof line(s) 4999; it names a standard fact of the same chapter that the counted argument uses, and every step of the counted proof built on it is retained verbatim in the delivered body.\par\medskip
\ifdefined\hypertarget\hypertarget{sourcepointer-2}{}\fi
\noindent\textit{Source reference.} Original definition 01VW: sources/stacks-chow/tags.txt, source lines 2461--2461. Prerequisite pointer invoked at proof line(s) 5007; it names a standard fact of a sibling chapter that the counted argument uses, and every step of the counted proof built on it is retained verbatim in the delivered body.\par\medskip
\ifdefined\hypertarget\hypertarget{sourcepointer-3}{}\fi
\noindent\textit{Source reference.} Original lemma 01VJ: sources/stacks-chow/tags.txt, source lines 2448--2448. Prerequisite pointer invoked at proof line(s) 5032; it names a standard fact of a sibling chapter that the counted argument uses, and every step of the counted proof built on it is retained verbatim in the delivered body.\par\medskip
\ifdefined\hypertarget\hypertarget{sourcepointer-4}{}\fi
\noindent\textit{Source reference.} Original lemma 0BCL: sources/stacks-chow/tags.txt, source lines 14284--14284. Prerequisite pointer invoked at proof line(s) 5006; it names a standard fact of a sibling chapter that the counted argument uses, and every step of the counted proof built on it is retained verbatim in the delivered body.\par\medskip
\ifdefined\hypertarget\hypertarget{sourcepointer-5}{}\fi
\noindent\textit{Source reference.} Original lemma 01W2: sources/stacks-chow/tags.txt, source lines 2467--2467. Prerequisite pointer invoked at proof line(s) 5023; it names a standard fact of a sibling chapter that the counted argument uses, and every step of the counted proof built on it is retained verbatim in the delivered body.\par\medskip
\ifdefined\hypertarget\hypertarget{sourcepointer-6}{}\fi
\noindent\textit{Source reference.} Original lemma 07RM: sources/stacks-chow/tags.txt, source lines 9872--9872. Prerequisite pointer invoked at proof line(s) 5038; it names a standard fact of a sibling chapter that the counted argument uses, and every step of the counted proof built on it is retained verbatim in the delivered body.\par\medskip
\begin{lemma}
\label{lemma-contract-when-quasi-projective}
Let $S$ be a Noetherian scheme. Let $f : X \to S$ be a morphism of finite type.
Let $E \subset X$ be an exceptional curve of the first kind which is in a
fibre of $f$.
\begin{enumerate}
\item If $X$ is projective over $S$, then there exists a contraction
$X \to X'$ of $E$ and $X'$ is projective over $S$.
\item If $X$ is quasi-projective over $S$, then there exists a contraction
$X \to X'$ of $E$ and $X'$ is quasi-projective over $S$.
\end{enumerate}
\end{lemma}

\begin{proof}
Both cases follow from Lemma \ref{lemma-contract-ample}
using standard results on ample invertible modules and
(quasi-)projective morphisms.

\medskip\noindent
Proof of (1). Projectivity of $f$ means that $f$ is proper and there exists
an $f$-ample invertible module $\mathcal{L}$, see
Morphisms, Lemma \ref{morphisms-lemma-projective-is-quasi-projective-proper}
and Definition \ref{morphisms-definition-quasi-projective}.
Let $U \subset S$ be an affine open containing the image of $E$.
By Lemma \ref{lemma-contract-ample} there exists a contraction
$c : f^{-1}(U) \to V'$ of $E$ and an ample invertible module
$\mathcal{N}'$ on $V'$ whose pullback to $f^{-1}(U)$ is equal to
$\mathcal{L}(nE)|_{f^{-1}(U)}$. Let $v \in V'$ be the closed point
such that $c$ is the blowing up of $v$.
Then we can glue $V'$ and $X \setminus E$ along
$f^{-1}(U) \setminus E = V' \setminus \{v\}$
to get a scheme $X'$ over $S$. The morphisms $c$ and
$\text{id}_{X \setminus E}$ glue to a morphism $b : X \to X'$
which is the contraction of $E$. The inverse image of $U$ in $X'$
is proper over $U$. On the other hand, the restriction of $X' \to S$
to the complement of the image of $v$ in $S$ is isomorphic to the
restriction of $X \to S$ to that open. Hence $X' \to S$ is proper
(as being proper is local on the base by
Morphisms, Lemma \ref{morphisms-lemma-proper-local-on-the-base}).
Finally, $\mathcal{N}'$ and $\mathcal{L}|_{X \setminus E}$ restrict to
isomorphic invertible modules over $f^{-1}(U) \setminus E = V' \setminus \{v\}$
and hence glue to an invertible module $\mathcal{L}'$ over $X'$.
The restriction of $\mathcal{L}'$ to the inverse image of $U$
in $X'$ is ample because this is true for $\mathcal{N}'$.
For affine opens of $S$ avoiding the image of $v$, we see that
the same is true because it holds for $\mathcal{L}$.
Thus $\mathcal{L}'$ is $(X' \to S)$-relatively ample by
Morphisms, Lemma \ref{morphisms-lemma-characterize-relatively-ample}
and (1) is proved.

\medskip\noindent
Proof of (2). We can write $X$ as an open subscheme of a scheme
$\overline{X}$ projective over $S$ by Morphisms, Lemma
\ref{morphisms-lemma-quasi-projective-open-projective}.
By (1) there is a contraction $b : \overline{X} \to \overline{X}'$
and $\overline{X}'$ is projective over $S$. Then we let
$X' \subset \overline{X}$ be the image of $X \to \overline{X}'$;
this is an open as $b$ is an isomorphism away from $E$.
Then $X \to X'$ is the desired contraction. Note that
$X'$ is quasi-projective over $S$ as it has an
$S$-relatively ample invertible module
by the construction in the proof of part (1).
\end{proof}

\end{document}
